\documentclass[12pt]{article}
\usepackage[T1]{fontenc}
\usepackage[utf8]{inputenc}
\usepackage{lmodern}
\usepackage{textcomp}
\usepackage{enumitem}
\usepackage{geometry}
\geometry{margin=1in}
\begin{document}
\begin{center}
Answer Key - Biology - Graduate Level Assessment 18 \
\small (Answers are ordered exactly as in the question paper.)
\end{center}

\section*{Multiple Choice}
\begin{enumerate}[label=\arabic*.]
\item \textbf{Answer:} A
\\ \textit{Options:} A) 50\%; B) 5\%; C) 60\%; D) 25\%; E) 90\%
\\ \textit{Note:} Approximately 50\% of the human genome encodes proteins because this fraction represents the portion of DNA that is transcribed into messenger RNA and ultimately translated into functional proteins, reflecting the complexity and diversity of protein-coding genes within the genome.
\item \textbf{Answer:} C
\\ \textit{Options:} A) The main source of variation in humans is genetic engineering.; B) The main source of variation in humans is due to dietary factors.; C) Shuffling of alleles is responsible for the greatest variation in humans.; D) The main source of variation in humans is due to exposure to radiation.; E) The main source of variation in humans is genetic drift.
\\ \textit{Note:} The shuffling of alleles during meiosis and sexual reproduction creates a diverse combination of genetic material in offspring, leading to the greatest variation among individuals in a population.
\item \textbf{Answer:} A
\\ \textit{Options:} A) The number of synapses crossed; B) The speed of the action potential; C) The amplitude of the action potential; D) The number of neurons activated; E) The refractory period of the neuron
\\ \textit{Note:} The correct answer is based on the principle that the greater the stimulus intensity, the more action potentials are generated, leading to an increased number of synapses crossed, which enhances the brain's ability to differentiate between varying intensities of stimuli.
\item \textbf{Answer:} E
\\ \textit{Options:} A) - .09176; B) .21185; C) .37748; D) .55962; E) .65361
\\ \textit{Note:} The correct answer of 0.65361 is derived from calculating the cumulative probabilities of a normal distribution for the values 5 and 11, using the mean of 9 and standard deviation of 3, to find the probability that X falls between these two values.
\item \textbf{Answer:} C
\\ \textit{Options:} A) parenchyma; B) phloem; C) endoderm; D) pericycle; E) cortex
\\ \textit{Note:} The endoderm is the innermost layer of the root cortex that regulates water and nutrient entry into the plant's vascular system, making it the primary point of entry for water absorbed by the roots of monocot plants.
\end{enumerate}

\section*{True / False}
\begin{enumerate}[label=\arabic*.]
\setcounter{enumi}{5}
\item \textbf{Answer:} False
\\ \textit{Options:} A) True; B) False
\\ \textit{Note:} Clinical variables, either isolated or as components of a model, could not identify all children with pathologic radiographs.
\item \textbf{Answer:} True
\\ \textit{Options:} A) True; B) False
\\ \textit{Note:} Both medial and lateral PTS were significantly steeper in failures of ACLR than the control group. Medial or lateral PTS \ensuremath{\geq} 5^\ensuremath{\circ} was a new risk factor of ACLR failure.
\item \textbf{Answer:} True
\\ \textit{Options:} A) True; B) False
\\ \textit{Note:} These findings show that phagocytic NADPH oxidase activity is increased in obesity and is related to preclinical atherosclerosis in this condition. We also suggest that hyperleptinemia may contribute to phagocytic NADPH oxidase overactivity in obesity.
\item \textbf{Answer:} True
\\ \textit{Options:} A) True; B) False
\\ \textit{Note:} Speed discrimination, per se, is not impaired in schizophrenia patients. The observed abnormality appears to be a consequence of impairment in generating or integrating the feedback information from eye movements. This study introduces a novel approach to motion perception studies and highlights the importance of concurrently measuring eye movements to understand interactions between these two systems; the results argue for a conceptual revision regarding motion perception abnormality in schizophrenia.
\item \textbf{Answer:} False
\\ \textit{Options:} A) True; B) False
\\ \textit{Note:} Patients with WD may possibly undergo cardiac surgery without a markedly enhanced risk for bleeding complications despite a more than usual transfusion requirement and significantly lower platelet counts perioperatively.
\end{enumerate}

\section*{Long Form Response}
\begin{enumerate}[label=\arabic*.]
\setcounter{enumi}{10}
\item \textbf{Answer:} The theoretical implications of genetic variation in a species with four alleles on chromosome 1 and three alleles on chromosome 2 indicate a significant level of potential genetic diversity. Specifically, the total number of distinct genotypes that can be formed is calculated by multiplying the number of alleles for each chromosome: 4 alleles on chromosome 1 and 3 on chromosome 2, resulting in 4 x 3 = 12 distinct genotypes. However, considering that each individual can be homozygous or heterozygous for these alleles, the total number of distinct genotypes expands to 60 possible combinations. Factors influencing this genetic diversity include mutation rates, selection pressures, genetic drift, and the mating system within the population, all of which can alter allele frequencies and contribute to the overall genetic structure of the species.
\\ \textit{Note:} correct because it accurately calculates the total number of distinct genotypes by multiplying the number of alleles on each chromosome, reflecting the genetic diversity possible in the species while acknowledging the influence of homozygosity and heterozygosity on genotype formation.
\item \textbf{Answer:} Cross-pollination in plants can be facilitated through a variety of methods, primarily categorized into natural processes and human intervention. While natural mechanisms such as wind, water, and pollinator activity (e.g., insects) play significant roles in promoting cross-pollination, the assertion that cross-pollination is solely dependent on human intervention is a critical perspective that merits analysis. In agricultural contexts, humans actively engage in methods such as controlled pollination, where specific plants are manually pollinated to ensure genetic diversity and enhance desirable traits, thereby circumventing the limitations of natural pollination processes. This method can lead to increased yields and improved crop resilience. However, reliance on human intervention could diminish the role of ecological interactions that naturally maintain plant diversity and adaptation, suggesting a complex interplay between anthropogenic practices and natural evolutionary processes in plant reproduction.
\\ \textit{Note:} correct because it comprehensively examines both natural and human-facilitated methods of cross-pollination, emphasizing the importance of understanding their distinct roles and implications in plant reproductive success.
\end{enumerate}

\vfill
\noindent\textit{End of Answer Key}
\end{document}